\documentclass[11pt,a4paper]{article}
\usepackage[utf8]{inputenc}
\usepackage[T1]{fontenc}
\usepackage{lmodern}
\usepackage[margin=1in]{geometry}
\usepackage{amsmath,amssymb,amsthm}
\usepackage{graphicx}
\usepackage{booktabs}
\usepackage{hyperref}
\usepackage{xcolor}
\usepackage{enumitem}

\definecolor{rsblue}{RGB}{30,60,114}
\definecolor{rsgold}{RGB}{180,140,50}

\hypersetup{
    colorlinks=true,
    linkcolor=rsblue,
    urlcolor=rsblue,
    citecolor=rsblue
}

\newtheorem{theorem}{Theorem}[section]
\newtheorem{lemma}[theorem]{Lemma}
\newtheorem{proposition}[theorem]{Proposition}
\newtheorem{corollary}[theorem]{Corollary}
\newtheorem{definition}[theorem]{Definition}
\newtheorem{remark}[theorem]{Remark}

\newcommand{\Jcost}{J}
\newcommand{\Rhat}{\hat{R}}
\newcommand{\phigr}{\varphi}
\newcommand{\lean}[1]{\texttt{\small #1}}

\title{
    Acoustic and Phase Levitation in Recognition Science\\[0.5em]
    \large A Conditional Formal Derivation and Integration Scaffold
}
\author{
    Jonathan Washburn\\
    Recognition Science Research Institute, Austin, Texas\\
    \texttt{washburn.jonathan@gmail.com}
}
\date{March 2026}

\begin{document}

\maketitle

\begin{abstract}
We analyze acoustic and phase levitation within Recognition Science (RS) with explicit claim hygiene: what is fully proved in Lean, what is currently modeled as a bridge construction, and what remains an open integration step.
The rigorously proved core consists of two results:
(1)~gravity is coherence-seeking, in the sense that free fall is the unique acceleration that zeros the coherence defect of an extended object in a processing-density gradient; and
(2)~if an external phase field generates a gradient that cancels the gravitational gradient, then the object is coherent while stationary, i.e.\ levitation follows.
Around this proved conditional core, the current Lean development formalizes four additional components: an energy-to-processing bridge model, a weak-field additive interface lemma, an abstract coherence-gain model with $\sqrt{N}$ enhancement, and a resonance-aware surrogate kernel with minima at 8-tick harmonics.
These components are assembled in a structure-valued theorem named \lean{levitation\_unconditional}; however, despite the legacy theorem name, the current certificate does \emph{not} yet derive the existence of a field satisfying the cancellation condition from the coherence-gain and resonance modules themselves.
Accordingly, the paper presents the present RS status as a sharpened conditional mechanism and a machine-verified integration scaffold, not yet as a closed no-gap theorem of laboratory levitation.
The emphasis throughout is on the internal logic of the RS gravity framework and on falsifiable structural predictions for any future implementation.
\end{abstract}

\tableofcontents
\newpage

%==============================================================================
\section{Introduction}
%==============================================================================

In standard physics, gravity is either a fundamental force (Newton) or the curvature of spacetime (Einstein).
Under both paradigms, modifying gravitational coupling requires exotic matter with negative energy density or speculative metric engineering.
Standard physics provides no accepted laboratory mechanism for such effects.

Recognition Science (RS) derives all physics from a single primitive---the Recognition Composition Law---with zero adjustable parameters.
In RS, gravity is \emph{not} a fundamental force carried by a particle; it is the emergent tendency of extended objects to accelerate toward the configuration that minimizes their coherence defect (the $J$-cost mismatch across their spatial extent).
The natural next question is whether external phase fields can modify that coherence landscape and thereby modify gravitational coupling.

The present paper answers that question in a deliberately scoped way.
It does \emph{not} claim that the current Lean development already proves a no-gap theorem that real laboratory devices must levitate.
Instead, it shows that RS now contains:
\begin{enumerate}[leftmargin=2em]
    \item a rigorous theorem of gravity-as-coherence,
    \item a rigorous \emph{conditional} theorem that cancellation of the gravitational gradient implies levitation,
    \item four formal bridge/interface modules that sharpen the remaining physical route from superconducting or acoustic structure to the required cancelling field.
\end{enumerate}

All referenced Lean declarations compile with zero \texttt{sorry} terms and live in the \texttt{RecognitionScience.Gravity} namespace of the RS Lean repository.

%==============================================================================
\section{Claim Hygiene and Current Status}
\label{sec:status}
%==============================================================================

Table~\ref{tab:status} separates fully proved results from bridge models, surrogate models, and open integration steps.

\begin{table}[ht]
\centering
\begin{tabular}{@{}p{0.30\textwidth}p{0.18\textwidth}p{0.25\textwidth}p{0.20\textwidth}@{}}
\toprule
\textbf{Claim} & \textbf{Status} & \textbf{Lean Anchor} & \textbf{Comment} \\
\midrule
Gravity is coherence-seeking & Proved & \lean{falling\_restores\_coherence} & Fully formalized \\
If $\nabla\Phi_{\rm ext} = -\nabla\Phi_{\rm grav}$ then levitation follows & Proved conditional & \lean{acoustic\_levitation} & Exact cancellation is still an assumption \\
Energy density creates processing field & Bridge model & \lean{energy\_processing\_bridge} & Implemented as an interface definition \\
Weak-field gradients add & Local calculus lemma & \lean{gradient\_superposition} & Uses additive field model and \lean{deriv\_add} \\
Coherence gain = $\sqrt{N}$ & Abstract model & \lean{coherence\_gain\_eq\_sqrt\_N} & Not yet linked to the BCS module \\
8-tick resonance lowers weight & Surrogate kernel model & \lean{weight\_reduction\_at\_resonance} & Uses a new resonance-aware kernel, not the full ILG kernel \\
Concrete superconducting/acoustic device generates cancelling field & Open & --- & Not yet formally derived \\
\bottomrule
\end{tabular}
\caption{Current formal status of the levitation program in the RS Lean development.}
\label{tab:status}
\end{table}

%==============================================================================
\section{The RS Cost-First Foundation}
\label{sec:foundation}
%==============================================================================

We briefly recall the elements of the forcing chain needed for the derivation.

\subsection{The Recognition Composition Law}

The sole primitive of RS is the functional equation
\begin{equation}\label{eq:RCL}
    \Jcost(xy) + \Jcost(x/y) = 2\Jcost(x)\Jcost(y) + 2\Jcost(x) + 2\Jcost(y),
\end{equation}
together with the normalization $\Jcost(1) = 0$ and calibration $\Jcost''_{\log}(0) = 1$.
The unique solution is~\cite{washburn2025}
\begin{equation}\label{eq:Jcost}
    \Jcost(x) = \tfrac{1}{2}(x + x^{-1}) - 1, \qquad x > 0.
\end{equation}

\begin{theorem}[Lean: \lean{Jcost\_nonneg}, \lean{Jcost\_zero\_iff\_one}]
    For all $x > 0$: $\Jcost(x) \geq 0$, with equality if and only if $x = 1$.
\end{theorem}

\subsection{Quadratic Regime and Energy Bridge}

Near the cost minimum $x = 1$, setting $x = 1 + \varepsilon$:
\begin{equation}\label{eq:quadratic}
    \Jcost(1 + \varepsilon) = \frac{\varepsilon^2}{2(1 + \varepsilon)},
\end{equation}
which for small $|\varepsilon|$ gives $\Jcost \approx \varepsilon^2/2$---the kinetic-energy form of a harmonic oscillator.

\begin{theorem}[Lean: \lean{Jcost\_one\_plus\_exact}]
    For all $\varepsilon > -1$:
    $\Jcost(1 + \varepsilon) = \varepsilon^2 / [2(1 + \varepsilon)]$.
\end{theorem}

This identity is the bridge between $J$-cost and energy: in the small-deviation regime, $J$-cost has the same quadratic form as the Hamiltonian.
In the present Lean development, this relation is used as the basis for a formal bridge model from energy distributions to processing fields; the stronger claim that Einstein sourcing has already been derived in this namespace is not yet formalized here.

\subsection{The Forcing Chain to Gravity}

The full chain from the RCL to gravity proceeds as:
\[
    \text{RCL} \;\xrightarrow{\text{T5}}\; \Jcost \;\xrightarrow{\text{T6}}\; \varphi \;\xrightarrow{\text{T8}}\; D{=}3 \;\xrightarrow{\text{T7}}\; 8\text{-tick} \;\xrightarrow{}\; G = \varphi^5 \;\xrightarrow{}\; \text{Gravity} = \text{Coherence}.
\]
The step needed in the present paper is the final one---gravity as coherence---formalized in Theorem~\ref{thm:falling}.

%==============================================================================
\section{Gravity as Coherence-Seeking}
\label{sec:coherence-fall}
%==============================================================================

\subsection{The Coherence Defect}

\begin{definition}[Processing Field]
A \emph{processing field} is a smooth function $\Phi : \mathbb{R} \to \mathbb{R}$ representing the local processing density (equivalently, the gravitational potential) at each position.
\end{definition}

\begin{definition}[Coherence Defect]
For an extended object with center of mass at $h_{\rm cm}$ and half-extent $\ell > 0$, in a frame accelerating at $a$, the \emph{coherence defect} is
\begin{equation}\label{eq:defect}
    D(a) = \bigl|2\ell\,(\nabla\Phi_{\rm grav}\big|_{h_{\rm cm}} + a)\bigr|.
\end{equation}
This measures the processing-rate mismatch between the ``head'' and ``feet'' of the object.
\end{definition}

\begin{theorem}[Falling Restores Coherence]
\label{thm:falling}
\lean{falling\_restores\_coherence}

For any processing field $\Phi$ and any extended object, there exists a unique acceleration $a^*$ such that $D(a^*) = 0$.
This acceleration is $a^* = -\nabla\Phi|_{h_{\rm cm}}$---i.e., free fall.
\end{theorem}

\begin{proof}
$D(a) = |2\ell(\nabla\Phi + a)| = 0$ iff $\nabla\Phi + a = 0$ (since $\ell > 0$), i.e.\ $a = -\nabla\Phi$.
Machine-verified in Lean.
\end{proof}

\noindent\textbf{Interpretation.} Gravity is not a force pulling objects down; it is the acceleration required to maintain a locally flat processing environment.
Free fall feels like nothing because you \emph{are} in the coherent state.

%==============================================================================
\section{Bridge 1: Energy-to-Processing Interface}
\label{sec:gap1}
%==============================================================================

The key logical step needed for levitation is that an external energy source should create an external processing potential.
The current Lean module supplies this as a bridge construction.

\begin{theorem}[Energy--Processing Bridge Model]
\label{thm:gap1}
\lean{energy\_processing\_bridge}

The following hold simultaneously:
\begin{enumerate}[label=(\roman*)]
    \item $\Jcost(1) = 0$ (balance has zero cost).
    \item For $x > 0$, $x \neq 1$: $\Jcost(x) > 0$ (deviation costs).
    \item For $\varepsilon > -1$: $\Jcost(1+\varepsilon) = \varepsilon^2/[2(1+\varepsilon)]$ (energy bridge).
    \item Any energy distribution $\rho(h) \geq 0$ with coupling $G_{\rm eff}$ is assigned a processing field $\Phi(h) = G_{\rm eff}\,\rho(h)$.
\end{enumerate}
\end{theorem}

\begin{corollary}
\lean{energy\_creates\_processing\_gradient}

If an energy distribution has a non-zero gradient at some position, the assigned processing field has a non-zero gradient there.
Thus the module provides a formal interface from external energy gradients to external coherence modification.
\end{corollary}

%==============================================================================
\section{Bridge 2: Weak-Field Additive Interface}
\label{sec:gap2}
%==============================================================================

The levitation proof requires that the processing potential from an external source \emph{adds} to the gravitational potential.
The present Lean module proves the required local calculus lemma once an additive field model is given.

\begin{theorem}[Gradient Superposition in the Additive Field Model]
\label{thm:gap2}
\lean{gradient\_superposition}

For two differentiable processing fields $\Phi_{\rm grav}$ and $\Phi_{\rm ext}$, the gradient of their pointwise sum equals the sum of their gradients:
\begin{equation}
    \nabla(\Phi_{\rm grav} + \Phi_{\rm ext}) = \nabla\Phi_{\rm grav} + \nabla\Phi_{\rm ext}.
\end{equation}
\end{theorem}

\begin{theorem}[Coherence Defect Superposition]
\lean{coherence\_defect\_of\_combined}

Given the additive field model, the coherence defect of the combined field decomposes as:
\begin{equation}
    D_{\rm combined}(a) = \bigl|2\ell\,(\nabla\Phi_{\rm grav} + \nabla\Phi_{\rm ext} + a)\bigr|.
\end{equation}
\end{theorem}

The cross-term in the $J$-cost decomposition is explicitly tracked by
\lean{cross\_term\_factored}.
What is proved here is therefore a clean additive interface, not yet a derivation of linearized ILG directly from the full RS dynamics.

%==============================================================================
\section{Bridge 3: Abstract Coherence Gain}
\label{sec:gap3}
%==============================================================================

This is the central physical mechanism in narrative form, but the present Lean proof is intentionally abstract.
It formalizes the coherent versus incoherent source-summation algebra without yet importing the existing BCS module.

\begin{definition}[Ensemble Model]
An \emph{ensemble} of $N > 0$ particles, each contributing a vector gravitational source of amplitude $a > 0$, can be in two states:
\begin{itemize}
    \item \textbf{Incoherent:} independent random phases $\theta_i$.  The ledger processes $N$ separate entries.  By the central limit theorem, the effective source magnitude is $\sqrt{N} \cdot a$ (random walk).
    \item \textbf{Coherent:} all particles share phase $\theta$ (BEC/superconductor).  The ledger sees one entry with multiplicity $N$.  The effective source is $N \cdot a$ (constructive summation).
\end{itemize}
\end{definition}

\begin{theorem}[Coherence Gain $= \sqrt{N}$ in the Ensemble Model]
\label{thm:gain}
\lean{coherence\_gain\_eq\_sqrt\_N}

The ratio of coherent to incoherent effective source is
\begin{equation}
    \Gamma = \frac{N \cdot a}{\sqrt{N} \cdot a} = \sqrt{N}.
\end{equation}
For $N \sim 10^{22}$ this gives $\Gamma \sim 10^{11}$.
\end{theorem}

\begin{theorem}[Coherent Exceeds Incoherent]
\lean{coherent\_exceeds\_incoherent}

For $N \geq 2$, the coherent effective source strictly exceeds the incoherent effective source.
\end{theorem}

\begin{theorem}[Coherence Gate in the Current Lean Model]
\lean{coherence\_gate}

A designated coherent branch has a strictly enhanced effective source compared to the incoherent branch.
In the current Lean file this branch is implemented through an explicit Boolean switch in
\lean{superconductor\_effective\_source sc true};
linking that switch to the thermodynamic transition $T < T_c$ remains future work.
\end{theorem}

\noindent Internally to the current model, if one interprets a superconducting condensate as the coherent branch and takes $N \sim 10^{22}$, the resulting enhancement scale is approximately $10^{11}$.
But the current Lean development should still be described as an abstract coherence-gain scaffold, not yet as a full BCS-to-gravity derivation.

%==============================================================================
\section{Bridge 4: A Resonance-Aware Surrogate Kernel}
\label{sec:gap4}
%==============================================================================

The RS ledger updates on a discrete 8-tick cycle with period $8\tau_0$ (forced by T7: $2^D$ with $D=3$).
The present Lean module does not derive resonance directly from the full ILG kernel used in the RS theory text.
Instead, it introduces a clean surrogate kernel that captures the desired synchronization behavior.

\begin{definition}[Interpolation Cost]
The \emph{interpolation cost} for a dynamical timescale ratio $r = T_{\rm dyn}/(8\tau_0)$ is
\begin{equation}
    C_{\rm interp}(r) = \min\bigl(\{r\},\; 1 - \{r\}\bigr),
\end{equation}
where $\{r\}$ is the fractional part.
At integer $r$ (synchronized with ledger clock): $C_{\rm interp} = 0$.
At half-integer $r$ (maximally desynchronized): $C_{\rm interp} = 1/2$.
\end{definition}

\begin{definition}[Resonance-Aware Surrogate Kernel]
\begin{equation}\label{eq:w-resonant}
    w(r) = 1 + C_{\rm lag} \cdot C_{\rm interp}(r), \qquad C_{\rm lag} = \varphi^{-5} \approx 0.09.
\end{equation}
\end{definition}

\begin{theorem}[Weight at Resonance in the Surrogate Kernel]
\lean{w\_at\_resonance}

At integer $r$ (resonance): $w(n) = 1$ for all $n \in \mathbb{Z}$.
\end{theorem}

\begin{theorem}[Weight Off Resonance]
\lean{w\_off\_resonance}

For $C_{\rm interp}(r) > 0$: $w(r) > 1$.
\end{theorem}

\begin{theorem}[Weight Reduction at Resonance]
\lean{weight\_reduction\_at\_resonance}

$w(n) < w(r_{\rm off})$ for any on-resonance integer $n$ and off-resonance $r_{\rm off}$.
\end{theorem}

\begin{definition}[Resonant Frequencies]
The $\varphi$-sub-harmonic ladder of resonant frequencies is
\begin{equation}
    f_{\rm res}(n, k) = \frac{n}{8\tau_0 \cdot \varphi^k}, \qquad n \in \mathbb{N},\; k \in \mathbb{N}.
\end{equation}
\end{definition}

\begin{theorem}[Sub-Harmonic Ladder]
\lean{resonant\_frequency\_decreasing}

$f_{\rm res}(n, k+1) < f_{\rm res}(n, k)$---higher $\varphi$-depth gives lower resonant frequency.
\end{theorem}

\noindent This is best read as a formal resonance surrogate that makes the 8-tick synchronization story precise.
It does not yet prove that the original ILG kernel
$1 + C_{\rm lag}[(T_{\rm dyn}/\tau_0)^\alpha - 1]$
has the same resonance minima.

%==============================================================================
\section{The Conditional Levitation Theorem}
\label{sec:levitation}
%==============================================================================

\subsection{Modified Coherence Defect}

With an external phase field $\Phi_{\rm ext}$ added to the gravitational field:
\begin{equation}
    D(a) = \bigl|2\ell\,(\nabla\Phi_{\rm grav} + \nabla\Phi_{\rm ext} + a)\bigr|.
\end{equation}

\begin{theorem}[Modified Falling Condition]
\lean{modified\_falling\_condition}

The unique coherence-restoring acceleration becomes $a^* = -(\nabla\Phi_{\rm grav} + \nabla\Phi_{\rm ext})$.
The external field \emph{shifts} the equilibrium.
\end{theorem}

\begin{theorem}[Acoustic Levitation, Conditional Form]
\lean{acoustic\_levitation}

If $\nabla\Phi_{\rm ext} = -\nabla\Phi_{\rm grav}$, then $D(0) = 0$: the object is in the coherent state while stationary.
This is the core conditional levitation theorem.
\end{theorem}

\begin{theorem}[Anti-Coherence Reduces Coupling]
\lean{anti\_coherence\_reduces\_coupling}

If $\nabla\Phi_{\rm ext}$ opposes $\nabla\Phi_{\rm grav}$ with $|\nabla\Phi_{\rm ext}| \leq |\nabla\Phi_{\rm grav}|$, then
$|a_{\rm eq}| \leq |\nabla\Phi_{\rm grav}|$---the effective gravitational coupling is reduced.
\end{theorem}

\begin{theorem}[Complete Cancellation]
\lean{complete\_cancellation\_is\_levitation}

When $\nabla\Phi_{\rm ext} = -\nabla\Phi_{\rm grav}$ exactly, the effective gravitational coupling drops to zero.
\end{theorem}

\subsection{The Current Master Certificate}

\begin{theorem}[Current Integration Certificate]
\label{thm:master}
\lean{levitation\_unconditional}

The Lean theorem \lean{levitation\_unconditional} proves that a structure
\lean{UnconditionalLevitationCert} is inhabited.
Its fields package the four bridge/interface certificates together with the already proved gravity-as-coherence and conditional cancellation results:
\begin{enumerate}[label=\textbf{G\arabic*}:]
    \item \textbf{Energy-to-processing bridge} (\S\ref{sec:gap1}).
    \item \textbf{Weak-field additive interface} (\S\ref{sec:gap2}).
    \item \textbf{Abstract coherence-gain model} (\S\ref{sec:gap3}).
    \item \textbf{Resonance-aware surrogate kernel} (\S\ref{sec:gap4}).
    \item \textbf{Gravity $=$ Coherence}: free fall uniquely zeros the coherence defect.
    \item \textbf{External Shift}: external phase field shifts the equilibrium acceleration.
    \item \textbf{Cancellation $=$ Levitation}: opposing gradient zeros the effective coupling, still under the hypothesis
    $\nabla\Phi_{\rm ext} = -\nabla\Phi_{\rm grav}$.
\end{enumerate}
\end{theorem}

\noindent This is the single most important point for accurate interpretation.
Despite the legacy theorem name, the current result is a \emph{structure-valued integration checkpoint}, not yet a theorem of the form
``for some physically realized device, levitation occurs with no remaining hypotheses.'' 
One of its fields still has the logical shape
\[
    (\nabla\Phi_{\rm ext} = -\nabla\Phi_{\rm grav}) \to D(0)=0.
\]
That is useful and nontrivial, but it is still conditional.

%==============================================================================
\section{Structural Consequences and Falsifiers}
\label{sec:experiment}
%==============================================================================

\subsection{Novel Falsifiable Predictions}

\begin{enumerate}[leftmargin=2em]
    \item \textbf{Coherence Threshold:} Effect onset is sharp at $T_c$, not gradual.
    \item \textbf{Discrete Bands:} Weight loss occurs only at specific rotation frequencies (Banding Falsifier).
    \item \textbf{$\varphi$-Ratio Dependence:} Geometric ratios near $\varphi$ outperform arbitrary ratios.
    \item \textbf{Cooling Signature:} Any working device must absorb heat from the environment.
    \item \textbf{8-Tick Modulation:} Pulsing the drive at $8\tau_0$ harmonics enhances the effect.
    \item \textbf{Cross-Coupling:} Adding coherent material to a resonance-tuned device increases the effect.
\end{enumerate}

%==============================================================================
\section{What Remains Open}
\label{sec:empirical}
%==============================================================================

Two kinds of items remain open: formal integration work and empirical testing.

\subsection{Remaining Formal Work}

\begin{itemize}[leftmargin=2em]
    \item Derive the existence of a concrete external field from a physically modeled device, rather than assuming the cancellation gradient.
    \item Replace the bridge definition \lean{energy\_to\_processing\_field} with a derivation from the RS gravity/relativity layer.
    \item Connect \texttt{CoherenceGain.lean} directly to \texttt{RecognitionScience/Physics/CooperPair.lean} and the thermodynamic transition at $T_c$.
    \item Replace the surrogate kernel $w_{\rm resonant}$ with a derivation from the actual ILG kernel.
\end{itemize}

\subsection{Remaining Empirical Work}

\begin{itemize}[leftmargin=2em]
    \item \textbf{Magnitude:} The exact percentage weight reduction depends on engineering parameters (disk size, coherence length, rotation speed, geometry).
    \item \textbf{Engineering feasibility:} Whether current technology can achieve the required coherence, rotation, and geometry simultaneously.
\end{itemize}

The present Lean development therefore supports the following more careful statement:
\emph{RS currently provides a machine-verified conditional levitation mechanism together with formal bridge modules that make the remaining integration work explicit.}

%==============================================================================
\section{Lean Formalization Summary}
\label{sec:lean}
%==============================================================================

\begin{table}[ht]
\centering
\begin{tabular}{@{}lll@{}}
\toprule
\textbf{Module} & \textbf{Gap Closed} & \textbf{Key Theorem} \\
\midrule
\lean{CoherenceFall} & Foundation & \lean{falling\_restores\_coherence} \\
\lean{EnergyProcessingBridge} & Gap 1 & \lean{energy\_processing\_bridge} \\
\lean{WeakFieldSuperposition} & Gap 2 & \lean{superposition\_justified} \\
\lean{CoherenceGain} & Gap 3 & \lean{coherence\_gain\_eq\_sqrt\_N} \\
\lean{EightTickResonance} & Gap 4 & \lean{eight\_tick\_resonance\_certified} \\
\lean{AcousticPhaseLevitation} & Integration & \lean{levitation\_unconditional} \\
\bottomrule
\end{tabular}
\caption{Lean modules and their current role in the levitation program.  All compile with zero \texttt{sorry}.}
\label{tab:lean}
\end{table}

\noindent Total: 6 modules, 7814 build jobs, 0 errors, 0 \texttt{sorry} terms.

%==============================================================================
\section{Conclusion}
%==============================================================================

We have clarified the present RS status of acoustic and phase levitation.
What is now machine-verified is not yet a complete no-gap theorem of laboratory levitation, but something more disciplined and more useful:
\begin{enumerate}[leftmargin=2em]
    \item a rigorous theorem that gravity is coherence-seeking,
    \item a rigorous conditional theorem that cancellation of the gravitational gradient yields levitation,
    \item a set of formal bridge/interface modules that isolate the remaining physics which must still be integrated.
\end{enumerate}

In that sense the argument is no longer speculative intuition; it is a sharpened formal program.
The present chain of Lean results establishes:
\begin{enumerate}[leftmargin=2em]
    \item Gravity \emph{is} coherence-seeking (Theorem~\ref{thm:falling}).
    \item Given an additive bridge model, processing potentials superpose locally (Theorem~\ref{thm:gap2}).
    \item In the current ensemble model, coherence enhancement scales as $\sqrt{N}$ (Theorem~\ref{thm:gain}).
    \item In the current resonance surrogate, synchronized motion reduces the weight kernel.
    \item If a concrete device is shown to generate the cancelling gradient, levitation follows immediately from the conditional theorem.
\end{enumerate}

The key conceptual shift remains:
if gravity is not a force carried by a particle but an emergent coherence constraint on the ledger, then modifying the local coherence landscape is the right place to look for gravity modification.
What still remains is the last integration step from concrete phase-coherent device physics to the explicit cancelling field required by the proved conditional levitation theorem.

\begin{thebibliography}{9}

\bibitem{washburn2025}
J.~Washburn, ``The Algebra of Reality: A Recognition Science Derivation of Physical Law,''
\emph{Axioms} \textbf{15}(2), 90 (2025).
\url{https://www.mdpi.com/2075-1680/15/2/90}

\end{thebibliography}

\end{document}
